\section{Implementation}
\label{sec:implementation}

\alaya{} is implemented in Rust as an embeddable library crate, prioritizing
zero-dependency deployment and deterministic behavior.

\subsection{Storage Layer}

All data resides in a single SQLite database using WAL (Write-Ahead Logging)
mode for concurrent read access during write operations. The schema comprises
seven tables: \texttt{episodes} (with FTS5 virtual table for full-text
search), \texttt{semantic\_nodes}, \texttt{impressions},
\texttt{preferences}, \texttt{embeddings}, \texttt{links} (the Hebbian
graph), \texttt{node\_strengths} (Bjork dual-strength values), and
\texttt{categories} (emergent ontology). Schema versioning uses SQLite's
\texttt{PRAGMA user\_version} with forward migration on database open.

The choice of SQLite over a dedicated vector database or graph database is
deliberate. The companion survey~\cite{hui2026taxonomy} reports that 70.5\%
of memory systems require an LLM for memory management operations and 42.0\%
require 1--2 external services. \alaya{} requires neither: all memory
management operations (consolidation, transformation, forgetting,
perfuming, ontology) execute locally against SQLite, with the LLM used only
during consolidation for knowledge extraction---and even that is pluggable
via the \texttt{ConsolidationProvider} trait.

\subsection{Provider Interface}

\alaya{} does not embed any LLM or embedding model. Instead, it defines a
\texttt{ConsolidationProvider} trait that applications implement:

\begin{lstlisting}[language=Rust, caption={Provider trait (simplified)}]
pub trait ConsolidationProvider {
    fn extract_knowledge(
        &self,
        episodes: &[Episode],
    ) -> Result<Vec<NewSemanticNode>>;
}
\end{lstlisting}

This design keeps \alaya{} independent of any specific model, API, or
inference framework. Applications can use GPT-4, Claude, Gemini, local
models, or rule-based extractors. The emergent ontology system operates
entirely without the provider---category discovery uses embedding similarity
and graph topology, with LLM labeling as an optional enhancement.

\subsection{Embedding Handling}

Embeddings are stored as little-endian \texttt{f32} BLOBs in SQLite.
Cosine similarity is computed in Rust without external linear algebra
libraries. The system operates in three embedding modes:

\begin{itemize}[nosep]
    \item \textbf{Full embeddings:} Applications provide embeddings with
    episodes and semantic nodes. Vector retrieval and embedding-based
    clustering are fully operational.
    \item \textbf{Partial embeddings:} Some nodes have embeddings, others do
    not. The system uses embeddings where available and falls back to BM25
    and graph signals for nodes without embeddings.
    \item \textbf{No embeddings:} The system operates on BM25 full-text
    search and graph traversal alone. Category discovery falls back to
    graph-only clustering (shared Hebbian link overlap).
\end{itemize}

\subsection{MCP Server}

An optional Model Context Protocol (MCP)~\cite{anthropic2024mcp} server
binary (\texttt{alaya-mcp}) exposes \alaya{}'s API as MCP tools, enabling
any MCP-compatible AI agent to use \alaya{} as its memory backend. The MCP
server is gated behind a Cargo feature flag (\texttt{mcp}) and adds
\texttt{rmcp}, \texttt{tokio}, and \texttt{schemars} as dependencies only
when enabled. The default library crate has four dependencies:
\texttt{rusqlite} (with bundled SQLite), \texttt{serde},
\texttt{serde\_json}, and \texttt{thiserror}.

\subsection{Testing}

\alaya{} follows test-driven development with 181 tests across three
categories: 154 unit tests covering individual store operations, graph
algorithms, lifecycle functions, and category CRUD; 9 integration tests
covering full lifecycle pipelines (episode ingestion through category
emergence); and 18 documentation tests ensuring code examples in the API
documentation compile and produce correct results. Continuous integration
runs the full test suite across three operating systems (Linux, macOS,
Windows) and two Rust toolchains (stable and MSRV 1.75).
